% Version 0.25. Exact native local-K interfaces, relative to classical finite local duality.
\section{Finite local duality on the specified local fields}
\label{sec:local-k-duality}

\status{This section supplies the three literal interfaces
\texttt{LocalFinitenessK}, \texttt{LocalTateDualityK} and
\texttt{LocalEulerCharK}.  The prime may be $2$, the coefficient modules
may be nonfree, and the coefficient prime equals the residue
characteristic.  Classical finite local duality and the local Euler
formula are accepted in their precise usual forms; all field,
cochain, coefficient, cyclotomic and cup-product identifications needed
to apply them to the source interfaces are written here.}

Fix a prime $\ell$, a finite intermediate field
\[
 \mathbb Q_\ell\subset K\subset\Omega=\overline{\mathbb Q}_\ell
       =\texttt{PadicAlgCl}(\ell),\qquad
 d=[K:\mathbb Q_\ell],\qquad G=\operatorname{Aut}_K(\Omega).
 \tag{LK.1}
\]
The symbol $\Omega$ denotes the source's specified algebraic closure,
not its metric completion.  For $r>0$, put
$m=\ell^r$ and $R=\mathbb Z/m\mathbb Z$.
Let $M$ be any finite $R$-module, with the discrete topology and a
continuous $R$-linear action $\rho$ of $G$.  Its source cohomology is
denoted by $H^i(G,M)$ below.

\paragraph{Classical input boundary.}
The local-field fact used is uniqueness and completeness of the
extended absolute value on a finite extension of $\mathbb Q_\ell$;
see Milne, \emph{Algebraic Number Theory}, v3.08,
\href{https://www.jmilne.org/math/CourseNotes/ANT.pdf}{Theorem~7.38}.
The finite local cohomology input is finiteness, vanishing above degree
two, and perfect cup-product duality for finite discrete modules over
a characteristic-zero nonarchimedean local field.  Its dual coefficient
is $A^D=\operatorname{Hom}_{\mathbb Z}(A,\Omega^\times)$, with evaluation
and the usual local Brauer invariant.  See Milne,
\href{https://jmilne.org/math/Books/ADTnot.pdf}{\emph{Arithmetic Duality
Theorems}}, second edition, I, Theorem~2.1 and Corollary~2.3,
pp.~26--28.  Finally we use the local Euler formula, I, Theorem~2.8,
p.~31.  These statements apply when the module order is invertible in
$K$, which holds here even for $\ell$-power order.  They do not require
the order to be prime to the residue characteristic.  The cited
\href{https://www.jmilne.org/math/Books/add/ADT2006.pdf}{author's errata}
were checked; the correction concerning higher local fields is not a
hypothesis of the one-dimensional local-field statements used here.

\begin{proposition}[LK1: the actual field, cyclotomic action and cochains]
\label{prop:lk-identifications}
The field and group in (LK.1) are a characteristic-zero nonarchimedean
local field and its absolute Galois group, on the original underlying
objects.  Reduction of the source's $\ell$-adic cyclotomic character is
the character $\chi=\texttt{cycK}(\ell,K,r)$ characterized by
\[
   g\zeta=\zeta^{\chi(g)}\qquad
   (g\in G,\ \zeta\in\mu_m(\Omega)).
   \tag{LK.2}
\]
For any primitive $m$th root $\zeta$, the map
$a\mapsto\zeta^a$ identifies the source representation
\texttt{charRep}$(\chi)$ with $\mu_m(\Omega)$, in additive notation.
The source's $H^0,H^1,H^2$ and $H^n$ are the usual continuous
cohomology groups, with their original $R$-module structures.
\end{proposition}

\begin{proof}
The norm on $\Omega$ extends the norm of $\mathbb Q_\ell$; its
restriction to the finite extension $K$ is therefore the unique
extension in the cited local-field theorem.  It is complete and
discrete.  Here is also the lattice consequence that will be needed
for the Euler factor.  Choose a $\mathbb Q_\ell$-basis $b_1,\ldots,b_d$
and put $\Lambda=\sum_i\mathbb Z_\ell b_i$.  The coordinate maximum
norm and the restricted norm are equivalent.  Indeed, the upper bound
follows from the ultrametric inequality; for the lower bound minimize
$\|\sum_i x_i b_i\|$ on the compact set
$\{(x_i)\in\mathbb Z_\ell^d:\max_i|x_i|_\ell=1\}$.
The minimum is positive because the basis is independent.
Scaling a nonzero coordinate vector to this set proves the lower bound.
Consequently, for some integer $a\ge0$,
\[
 \ell^a\Lambda\subset\mathcal O_K\subset\ell^{-a}\Lambda.
 \tag{LK.3}
\]
The valuation ring $\mathcal O_K$ is a submodule of a finite free
$\mathbb Z_\ell$-module and contains a full lattice.  Since
$\mathbb Z_\ell$ is a principal ideal domain, it is free of rank $d$.
In particular its residue field is finite: it is a quotient of the
finite group $\mathcal O_K/\ell\mathcal O_K$.
Thus $K$ satisfies precisely the local-field hypotheses of the
classical theorems.  No completion changes its elements or its degree.

Every element of $\Omega$ is algebraic and separable over $K$, and
$\Omega$ is algebraically closed, so it is a separable closure of
this very $K$.  The source group $G$ is exactly
$\operatorname{Gal}(\Omega/K)$ with its Krull topology.  The subgroups
fixing finite extensions of $K$ give its neighborhood basis.  If one
describes the topology through the action on the valued algebraic
closure, the same basis results: each algebraic element has finitely
many conjugates, and a sufficiently small ball separates its chosen
value from the other conjugates.  This is also why every algebraic
$K$-automorphism is compatible with the local-field topology; the
unique extended norm is invariant under it.

Choose compatible primitive roots $\zeta_s$ of order $\ell^s$ with
$\zeta_{s+1}^{\ell}=\zeta_s$.  Write
$g\zeta_s=\zeta_s^{a_s(g)}$.  The residues $a_s(g)$ are units and
are compatible under reduction.  Their inverse limit is the
$\ell$-adic cyclotomic character.  A different compatible primitive
system multiplies its generators by a compatible system of units and
leaves every $a_s(g)$ unchanged.  Its reduction modulo $\ell^r$ is
therefore exactly the source definition of $\chi$, proving (LK.2)
for every $m$th root.  Its kernel fixes the finite extension $K(\mu_m)$,
so $\chi$ is continuous.  A primitive root exists in the specified
algebraic closure and its powers give a bijection $R\simeq\mu_m$;
(LK.2) proves equivariance.  Both finite coefficient groups are
discrete.  This does not assert that the roots lie in $K$.

For the given action, continuity means that every orbit map is
continuous.  Because $M$ is finite and discrete, the intersection of
the stabilizers of its elements is an open normal subgroup.  Thus the
action factors through a finite quotient, as required for a finite
discrete Galois module.

To check cohomology itself, the source uses the continuous
inhomogeneous cochains $C^n(G,M)=\operatorname{Cont}(G^n,M)$ and
the differential
\[
\begin{split}
 (\partial f)(g_1,\ldots,g_{n+1})={}&
 g_1f(g_2,\ldots,g_{n+1})\\
 &+\sum_{i=1}^{n}(-1)^i
 f(g_1,\ldots,g_i g_{i+1},\ldots,g_{n+1})\\
 &+(-1)^{n+1}f(g_1,\ldots,g_n).
\end{split}
\tag{LK.4}
\]
In degree zero this gives the invariant submodule.  In degrees one
and two the source spells out respectively
\[
 c(gh)=c(g)+g c(h),\qquad
 g z(h,k)-z(gh,k)+z(g,hk)-z(g,h)=0;
 \tag{LK.5}
\]
its boundaries are $g\mapsto gm-m$ and
$(g,h)\mapsto g b(h)-b(gh)+b(g)$, with $b$ continuous.
These are exactly (LK.4).  In every higher degree the source's
\texttt{Fin.contractNth} terms multiply adjacent arguments, and the
last term deletes the final argument, giving the displayed formula.
Thus kernels and quotients agree, rather than merely having equal
cardinalities.  Forgetting $R$ to the underlying abelian group changes
neither the cochains nor these differentials.  Multiplication by
every element of $R$ is represented by multiplication by an integer,
so the identification preserves the $R$-module structures as well.
\end{proof}

\begin{proposition}[LK2: literal local finiteness]
\label{prop:lk-finiteness}
For all the original data above, $H^0(G,M)$, $H^1(G,M)$ and
$H^2(G,M)$ are finite, and $H^n(G,M)=0$ for $n\ge3$.
Consequently the source's \texttt{LocalFinitenessK}
holds with its full quantifiers.
\end{proposition}

\begin{proof}
LK1 gives the actual finite local field, its specified separable
closure and a finite discrete continuous module.  Its order is a
positive integer and hence invertible in the characteristic-zero
field $K$.  The finite local cohomology theorem therefore gives
finiteness and vanishing for the continuous cohomology of these
objects.  LK1 identifies these groups with the source groups in each
degree, including the separate definition \texttt{Hn} in degrees
at least three.  A zero cohomology group has exactly one element,
which is the source's \texttt{Subsingleton} conclusion.  Neither
freeness over $R$ nor triviality of the action has been used.
\end{proof}

\begin{proposition}[LK3: invariant and all four literal cup bijections]
\label{prop:lk-tate}
Put $P=H^2(G,\texttt{charRep}(\chi))$ and
$M^D=\operatorname{Hom}_R(M,R)$ with action
\[
 (g\varphi)(x)=\chi(g)\varphi(g^{-1}x).
 \tag{LK.6}
\]
There is an $R$-linear isomorphism $j:P\simeq R$.  Each of the
following source maps is bijective:
\[
\begin{aligned}
 H^0(G,M)&\longrightarrow\operatorname{Hom}_R(H^2(G,M^D),P),\\
 H^1(G,M^D)&\longrightarrow\operatorname{Hom}_R(H^1(G,M),P),\\
 H^1(G,M)&\longrightarrow\operatorname{Hom}_R(H^1(G,M^D),P),\\
 H^2(G,M)&\longrightarrow\operatorname{Hom}_R(H^0(G,M^D),P).
\end{aligned}
\tag{LK.7}
\]
They are, in order, \texttt{cupH0H2TwistedDual},
\texttt{cupH1H1TwistedDual}, its \texttt{flip}, and
\texttt{cup20}, all with the source's identity group homomorphism.
Thus the literal \texttt{LocalTateDualityK} holds.
\end{proposition}

\begin{proof}
First consider coefficients carefully.  Every additive homomorphism
from $M$ to $\Omega^\times$ has image in $\mu_m$, because $mM=0$.
Via $a\mapsto\zeta^a$ from LK1 it becomes a unique map $M\to R$.
An additive map between $R$-modules is $R$-linear, since scalar
multiplication is integer multiplication modulo $m$.  Therefore
the classical Tate dual is the source's \texttt{TDual} with action
(LK.6), exactly \texttt{twDual}; the inverse on the argument is
essential.  Its action is continuous: the common open kernel of
$\rho$ and $\chi$ acts trivially on this finite dual.

No coefficient freeness is implicit.  Write a finite $R$-module as
a direct sum of $R/(\ell^{a_i})$, $1\le a_i\le r$; the zero module
corresponds to the empty sum.  On such a summand,
\[
 \operatorname{Hom}_R(R/(\ell^a),R)=\ell^{r-a}R.
 \tag{LK.8}
\]
The evaluation map into the double dual is injective: a nonzero
residue $x$ modulo $\ell^a$ is detected by the functional taking
$1$ to $\ell^{r-a}$.  Both modules have order $\ell^a$, so the
evaluation is an isomorphism.  Taking direct sums proves this for
all finite $R$-modules.  For the twisted duals it is equivariant,
since the two factors $\chi(g)$ and $\chi(g)^{-1}$ cancel in
evaluation.

The Kummer sequence and Hilbert~90 identify
$H^2(G,\mu_m)$ with the $m$-torsion of $\operatorname{Br}(K)$.
The local Brauer invariant identifies that torsion with
$m^{-1}\mathbb Z/\mathbb Z$.  Composing with the inverse of
\[
 \iota_m:R\longrightarrow m^{-1}\mathbb Z/\mathbb Z,
 \qquad a\longmapsto a/m,
 \tag{LK.9}
\]
and the LK1 coefficient identification gives $j:P\simeq R$.
This can be chosen before $M$ and $\rho$.  All its maps are additive
and hence $R$-linear.  The opposite common invariant sign, as in
LD's arithmetic-Frobenius normalization, multiplies $j$ by $-1$ and
leaves all the bijectivity conclusions unchanged.  The local-K
interface asks for an isomorphism, not an extra normalization.

For clarity the actual source cochain pairings in (LK.7) are as
follows.  If $x\in M^G$, $z_D$ is a dual-valued two-cocycle,
$f$ is a dual-valued one-cocycle, $c$ is an $M$-valued one-cocycle,
$z$ is an $M$-valued two-cocycle, and $\varphi\in(M^D)^G$, they are
\[
\begin{aligned}
  (x\cup z_D)(g,h)&=z_D(g,h)(x),\\
  (f\cup c)(g,h)&=f(g)\bigl(\rho(g)c(h)\bigr),\\
  (z\cup\varphi)(g,h)&=\varphi\bigl(z(g,h)\bigr).
\end{aligned}
\tag{LK.10}
\]
These follow directly from the source's evaluation pairing and its
continuous inhomogeneous cup products.  In the last formula the
usual action of $gh$ on the degree-zero second argument disappears
because $\varphi$ is invariant.  Thus these are the classical cup
products before the invariant, with the displayed order of the
arguments; the degree-one source map puts the dual first.

Finite local Tate duality, applied to $M$ and to $M^D$ with the
evaluation identification $M^{DD}\simeq M$, now says that composing
the first, second and fourth maps with the invariant gives perfect
pairings into $\mathbb Q/\mathbb Z$.  All cohomology groups here
are finite by LK2 and are killed by $m$, so every additive character
on them lands in $m^{-1}\mathbb Z/\mathbb Z$.  Equation (LK.9)
therefore identifies their character groups with their $R$-linear
duals.  This proves bijectivity after composing the target with $j$.
Postcomposition
\[
 \operatorname{Hom}_R(Y,P)\longrightarrow\operatorname{Hom}_R(Y,R),
 \qquad u\longmapsto j\circ u
 \tag{LK.11}
\]
is itself an isomorphism, with inverse $v\mapsto j^{-1}\circ v$.
Consequently the original maps into $P$, not only the scalar-valued
maps, are bijective.

Finally let $X=H^1(G,M^D)$ and $Y=H^1(G,M)$, and write
$b:X\to Y^*$ for the already proved scalar-valued isomorphism.
The flipped pairing is the composition
\[
 Y\xrightarrow{\ \mathrm{ev}\ }Y^{**}
   \xrightarrow{\ b^*\ }X^*,
 \qquad y\longmapsto\bigl(x\longmapsto b(x)(y)\bigr).
 \tag{LK.12}
\]
Both arrows are isomorphisms by finite biduality and bijectivity of
$b$.  Applying (LK.11) once more proves the exact \texttt{flip}
clause.  This uses the same bilinear form with its arguments exchanged;
no alternatingness or division by $2$ is involved.  It is valid also
when $\ell=2$.
\end{proof}

\begin{proposition}[LK4: literal residue-characteristic Euler factor]
\label{prop:lk-euler}
For the same original field and arbitrary finite coefficient module,
\[
 |H^1(G,M)|=|H^0(G,M)|\,|H^2(G,M)|\,|M|^{[K:\mathbb Q_\ell]}.
 \tag{LK.13}
\]
Hence the source's \texttt{LocalEulerCharK} holds.
\end{proposition}

\begin{proof}
Put $q=|M|$.  From the finite module decomposition in (LK.8),
$q=\ell^s$ for some $s\ge0$.  LK1 proves that $\mathcal O_K$ is
free of rank $d$ over $\mathbb Z_\ell$.  Reduction of a basis gives
\[
 |\mathcal O_K/q\mathcal O_K|
   =|\mathbb Z_\ell/\ell^s\mathbb Z_\ell|^d
   =\ell^{sd}=q^d.
 \tag{LK.14}
\]
This computation uses the rank of the local field's ring of
integers; it does not identify $M$ with a free coefficient module.
It includes $s=0$, when the displayed quotient is the zero group
and has order one.

The classical local Euler formula for this finite discrete module is
\[
 \frac{|H^0(G,M)|\,|H^2(G,M)|}{|H^1(G,M)|}
     =|\mathcal O_K/q\mathcal O_K|^{-1}.
 \tag{LK.15}
\]
All three cohomology groups are finite by LK2, and each contains zero,
so their cardinalities are positive.  Substitute (LK.14) and clear
the positive denominators to obtain (LK.13).  LK1's identifications
preserve all cardinalities, and its $d$ is exactly the source's
\texttt{Module.finrank}$(\mathbb Q_\ell,K)$.  Thus (LK.13) is the
literal equality of natural cardinalities in \texttt{LocalEulerCharK},
with the residue-characteristic factor retained.
\end{proof}

\paragraph{Closed interface and remaining scope.}
LK2, LK3 and LK4 can now be inserted directly as the three local-K
inputs wherever those literal predicates occur, including the
\texttt{LemmaK} and two-primary/reference-layer chains.  Their
quantifiers cover every finite intermediate $K$, every $r>0$, every
finite discrete $R$-module and every continuous action.  This does
not by itself prove the global integral-cohomology reduction, a
$K$-theoretic comparison, an arithmetic-group torsion bound, or the
odd-plus count/vanishing suppliers.  Those remain their own
mathematical chains.
